% chapters/04/rangkuman.tex
\bagian{Rangkuman}

\begin{rangkumanbox}
\begin{itemize}
  \item Siklus permainan terdiri atas tiga tahap---masukan, pembaruan, dan
        penayangan---yang dikerjakan berurutan dan diulang terus-menerus
        selama permainan berjalan.
  \item Urutan tahapan bersifat wajib: masukan lebih dahulu agar pembaruan
        mengetahui kehendak pemain, dan penayangan paling akhir agar yang
        terlihat adalah keadaan yang sudah selesai dihitung.
  \item Variabel keadaan harus dibuat di luar perulangan agar nilainya bertahan
        dari satu bingkai ke bingkai berikutnya.
  \item \code{pygame.event.get()} mengambil seluruh isi antrean peristiwa tanpa
        menunggu, sehingga siklus tetap berjalan walaupun pemain sedang diam.
        \code{pygame.event.wait()} menghentikan program dan tidak boleh dipakai
        di dalam siklus.
  \item Antrean peristiwa harus dikosongkan setiap bingkai; membiarkannya
        menumpuk membuat jendela tidak menanggapi masukan.
  \item Peristiwa berbeda dari keadaan tombol. Peristiwa \code{KEYDOWN} terjadi
        sekali saat tombol mulai ditekan, sedangkan gerak yang berlangsung
        selama tombol ditahan memerlukan \code{pygame.key.get_pressed()}.
  \item \code{clock.tick()} mengukur selang waktu satu bingkai sekaligus
        menahan laju bingkai, sehingga perintah itu diletakkan pada akhir
        putaran siklus.
  \item Batas laju bukan pembatasan mutu, melainkan penghematan. Pengukuran
        memperlihatkan bahwa siklus berbatas memakai beberapa persen prosesor,
        sedangkan siklus tanpa batas memakai hampir seluruh satu inti
        prosesor.
  \item Menaikkan batas laju tidak mempercepat program yang pekerjaannya memang
        berat; yang menentukan laju tercapai adalah lama pekerjaan satu
        bingkai.
  \item Siklus pada bab ini memakai langkah berubah. Nilai selang waktu sudah
        diukur, tetapi baru dipakai untuk menghitung waktu berjalan; Bab~V
        memakainya untuk menggerakkan objek.
\end{itemize}
\end{rangkumanbox}
